% \chapter{Introduction}
% % PREAMBLE (max 4 pages / ~16 paragraphs)

% \section{Contributions}

% \section{Thesis Outline}

% \section{Published Work}

% % Short section describing where the work presented in each chapter was published (with citation) 

\chapter{Introduction}
\label{chap:introduction}


Reinforcement learning (RL) offers a general framework for learning behaviour from interaction. Instead of manually specifying a controller, an objective is specified through rewards, and an agent improves its policy by acting, observing the consequences of its actions, and updating its behaviour accordingly \citep{sutton2018reinforcement}. This makes RL a powerful paradigm for adaptive control: RL algorithms can, in principle, discover behaviours that are difficult to design by hand, exploit large-scale computation during training, and deploy policies whose inference cost is comparatively low.

For many real-world systems, however, unrestricted interaction during training is impractical. Physical interaction can be slow, costly, unsafe, or destructive. Collecting millions of interactions with a real robot may be prohibitively time-consuming and economically unviable; an autonomous vehicle cannot explore dangerous manoeuvres on public roads; and a controller for critical infrastructure cannot learn through arbitrary trial and error. Online RL training is therefore commonly transferred from the deployment environment to a simulator, where experience can be generated more quickly and safely. This transfer removes many of the immediate costs of physical interaction. It does not, however, make suitable training experience available automatically. Instead, progress can become constrained by the ability to construct the simulated environments from which that experience is obtained. This thesis calls this constraint the \emph{simulation bottleneck}.

%This dependence introduces what this thesis terms the \emph{simulation bottleneck}. Effective simulator-based RL generally presupposes two conditions. First, there must exist a sufficiently accurate and interactable simulator of the target environment. Second, the simulator must expose the agent to a sufficiently diverse set of scenarios to support generalisation to the variations encountered after deployment. Neither condition is trivial. Constructing a simulator can require substantial domain expertise, software engineering, system identification, and validation. Even when a simulator already exists, designing a distribution of training scenarios that captures the relevant variability of the deployment environment can itself become a major engineering task.

\begin{tcolorbox}[
colback=gray!10,
leftrule=0.5mm,
top=3mm,
bottom=3mm,
title={The simulation bottleneck}
]
The \emph{simulation bottleneck} is the engineering burden incurred when the interaction required by RL must be produced through simulation rather than obtained freely from the deployment environment. Effective simulator-based RL generally requires two conditions:
\begin{enumerate}[leftmargin=*]
\item \textbf{Simulation capability.}
A simulator must be available that responds to the agent's actions and reproduces, with sufficient accuracy, the task-relevant observations, dynamics, and constraints of the target environment.

\item \textbf{Distributional coverage.} The simulator must expose the agent to a sufficiently diverse and deployment-relevant distribution of environment instances. These instances must collectively cover the variations in initial states, layouts, objects, dynamics, appearances, and other factors to which the learned policy will be expected to generalise at deployment.
\end{enumerate}

Neither condition follows from the design of the RL algorithm or model architecture. Satisfying the first often requires substantial domain expertise, system identification, software engineering, and validation. Satisfying the second requires identifying the relevant sources of environmental variation and constructing a training distribution that represents them adequately. If either condition is not met, agent performance is constrained by the process through which its experience is generated, rather than solely by the capabilities of the learning algorithm. Together, these conditions form a bottleneck that has traditionally required substantial human engineering to overcome.
\end{tcolorbox}

The simulation bottleneck exposes a mismatch between the promise of RL as a data-driven method and the human engineering often required to make it effective. RL is intended to reduce the need to specify behaviour manually, yet deploying an RL system may instead require much of the world in which that behaviour is learned to be specified manually. The complexity removed from policy design can therefore reappear in simulator construction and training-scenario generation.

The second condition is closely tied to the problem of generalisation. In practice, an agent is trained on only a finite collection of environment \textit{instances}, each representing a particular realisation of the underlying task. Strong performance on these training instances does not guarantee comparable performance on unseen instances, even when they are drawn from the same distribution \citep{procgen}. This phenomenon, commonly described as \emph{instance overfitting} \citep{instance_invariant}, undermines the expectation that a learned policy will transfer beyond the specific experience encountered during training.

This thesis investigates the simulation bottleneck through three connected questions. First, how do the representations learned by deep RL agents mediate instance overfitting? If agents trained on restricted sets of environment instances fail to generalise, one must understand whether this failure is visible in their internal representations, and whether architectural choices can reduce harmful instance-specific encoding. The first part of the thesis addresses this question by studying the interplay between actor and critic representations in modern actor--critic algorithms. It introduces mutual-information-based quantities for measuring representational dependence on environment identity, value structure, and latent dynamics, and relates these quantities to deployment generalisation error.

Second, if limited scenario diversity is a major cause of overfitting, can the set of training instances be expanded without introducing distributional shift with respect to the underlying task distribution? Resampling or prioritising existing instances can improve how a fixed training set is used, but it does not expand the number of available training scenarios. Conversely, fully unconstrained environment generation can create novel instances, but may drift away from the distribution that matters at deployment. The second part of this thesis addresses this tension through Data-Regularised Environment Design (\hbox{DRED}), which uses a learned generative model to augment a limited set of environment instances with synthetic but distributionally grounded scenarios.

% [TODO: Add a connection to the sim2real gap.]
% - Define it: gap between simulated and real data
% - Manually Engineering a simulator that fully bridge the gap is extrmeely costly/impossible
% - Learning a simulator directly from data is a promising way of bridging the gap / grounding simulators in reality

Third, can the simulator itself be learned from data? Even if scenario generation reduces the burden of manually designing a diverse training distribution, it still assumes that a suitable simulator already exists. This leaves the strongest form of the simulation bottleneck intact, and introduces the related \emph{sim-to-real gap} \citep{jakobi1995noise,tobin2017domain}: the discrepancy between the observations and dynamics represented in simulation and those encountered in the real world. Bridging this gap through manual simulator engineering requires accurately modelling complex physical interactions, visual variability, and sources of uncertainty. Doing so can be prohibitively expensive and, for sufficiently complex environments, may be practically impossible.

Learning a simulator directly from real interaction data offers a promising alternative, since it grounds the simulated experience in observations collected from the target environment. Here, we use the term \textit{world model} \citep{haworldmodels2018} to denote a learned, action-conditioned simulator: a model that predicts how the environment evolves under candidate actions and can therefore generate trajectories beyond those contained in the offline dataset. This allows agents to train or plan using simulated interaction learned from real-world data rather than a manually engineered simulator \citep{hafner2025training}. The final part of this thesis develops \hbox{PERSIST}, a world model that maintains a persistent latent 3D state in order to improve spatial memory, geometric consistency, and long-horizon stability in generated interactive experience.

Together, these questions define a progression. The thesis first studies what happens when RL agents learn from restricted simulator distributions; then develops a method for expanding such distributions using data-grounded generation; and finally investigates how interactable simulation can itself be learned from recorded experience.

\section{Contributions}
\label{sec:intro_contributions}
The principal contributions of this thesis are as follows.

\paragraph{Quantifying the information content of learned representations.}
We introduce mutual-information-based metrics for quantifying the information extracted by the learned representations of actor--critic agents. Specifically, these metrics measure the information encoded about inverse dynamics, forward dynamics, state value, and the identity of the environment instance from which an observation originates. We use them to establish an upper bound on an agent's deployment generalisation error that depends on the mutual information between its representation and the training-instance identity. We then conduct a systematic study of shared and separate actor--critic architectures, demonstrating that decoupled representations specialise in different forms of task-relevant information: actor representations preferentially retain information supporting action selection, while critic representations encode more information about value and environmental dynamics. This specialisation explains why representational decoupling can improve both capability and generalisation despite using parameters more efficiently: it limits interference between the actor and critic objectives and prevents information useful for fitting the critic on the training set from inducing instance-specific overfitting in the actor representation. Finally, we use these metrics to expose a common mechanism through which representation-learning objectives and adaptive instance-sampling strategies reduce overfitting: they promote the extraction of instance-agnostic information at the expense of information identifying individual training instances.

\paragraph{The role of the critic in data collection and exploration.}
We extend the mutual-information analysis from learned representations to the raw observations collected by the policy during training, revealing that the critic influences which data the agent encounters. This effect persists even when the actor and critic networks have fully separate parameters. Because the critic's value estimates enter the policy-gradient objective, they shape the actor's updates and thereby alter the policy's state-visitation distribution. Our results suggest that the policy is partly directed towards collecting experience that facilitates optimisation of the critic's own learning objective. This mechanism identifies a previously under-examined role for the critic in exploration and data collection. It may be beneficial when the critic steers the policy towards informative experience, but it may also constitute a residual source of objective interference, through which the critic biases policy learning even after their representations have been decoupled.

\paragraph{Training-instance augmentation through a learned generative model.}
We introduce Data-Regularised Environment Design (\hbox{DRED}), a method for improving zero-shot generalisation when only a limited set of training instances is initially available. DRED trains a generative model to approximate the distribution represented by the available environment parameters, and uses an adaptive sampling strategy to draw generated instances from this learned distribution during agent training. In this way, it expands the support of the training set rather than merely changing the frequency with which a fixed collection of instances is revisited. We show that DRED reduces the mutual information between the agent's representation and training-instance identity more effectively than methods that sample only from a fixed set of instances, including adaptive sampling methods. Consistent with the generalisation bound developed in this thesis, this reduction in instance-specific representation content is accompanied by lower generalisation error on previously unseen environments.

\paragraph{The importance of distributional grounding for zero-shot transfer.}
We demonstrate that existing unsupervised environment-design methods can induce substantial distribution shift between the environments generated during training and the target distribution used for evaluation. As a result, generating environments that are novel or difficult does not necessarily improve transfer and can instead reduce performance on the instances that matter at deployment. DRED prevents this failure by constraining environment generation through a learned approximation to the ground-truth instance distribution. Across the evaluated domains, this distributionally grounded augmentation improves zero-shot transfer over prior unsupervised environment-design approaches while preserving the relevance of generated instances to the target task. DRED therefore enables agents to extract stronger zero-shot generalisation capabilities from a restricted initial set of training instances, reducing the manual scenario-design burden associated with the second condition of the simulation bottleneck.

\paragraph{Neural simulators with persistent 3D state.}
We introduce \hbox{PERSIST}, an interactive world-modelling framework that maintains and evolves a persistent latent 3D representation of the environment. PERSIST structures simulation around three components---environment state, camera dynamics, and rendering---and projects the evolving 3D state into the agent's current egocentric view to condition visual generation. Unlike approaches that rely on a bounded history of pixel observations as their sole memory, this design retains spatial information about occluded and previously visited regions beyond the visual context window. Quantitative evaluations and a qualitative user study show that the resulting model improves spatial memory, geometric consistency, and long-horizon stability, sustaining coherent evolving worlds over extended interactions. The explicit 3D representation also enables forms of structured control that are difficult to obtain from purely pixel-based models, including synthesising diverse 3D environments from a single image and editing or specifying generated environments directly in 3D space. By advancing persistent, interactable neural simulators learned from offline interaction data, PERSIST addresses the first condition of the simulation bottleneck: the requirement that a suitable simulator be manually engineered before RL training can begin. Learning the simulator from observations of the target domain also provides a route towards narrowing the sim-to-real gap, since its generated experience is grounded in recorded environmental dynamics and visual observations rather than being specified entirely through manually constructed models.

\paragraph{A hierarchy of interventions on the simulation bottleneck.}
Taken together, these contributions establish a hierarchy of interventions on the simulation bottleneck. At the level of the agent, we identify and mitigate representational mechanisms that cause overfitting to a restricted set of simulated environments, while revealing how the critic continues to influence data collection even after architectural decoupling. At the level of the training distribution, we use generative modelling to produce additional scenarios while explicitly controlling distribution shift. At the level of the simulator, we introduce a neural simulation framework that can be learned from interaction data and maintain coherent experience over extended temporal horizons. The thesis therefore progresses from understanding the consequences of restricted simulation, to expanding the scenarios available within a simulator, and finally towards learning the simulator itself.

% \paragraph{Quantifying information content of learned representations.}
%[WRITE BASED ON FOLLOWING STRUCTURE BELOW]
% - we propose mutual information metrics to quantify information that actor-critic models learned representations extract from state or observations. Specifically: inverse and forward dynamics, environment instance specific information and state value information
% - We employ these metrics to:
%       - derive an upper bound on the generalisation error of RL algorithms that is proportional to [I(Z;L)]
%       - demonstrate information specialisation when actor and critic decouple their learned representations via a systematic study of shared and separate actor--critic architectures (chapter 2). Explains why representation decoupling improves agents capability & generalisation and is actually more parameter efficient -> prevents gradient (objective) interference and reduces overfitting of the actor representation to the training set.
%       - reveal the underlying mechanism of how certain representation learning (chapter 2) and adaptive instance sampling strategies (chapter 3) help prevent overfitting by promoting a representation learning to capture instance-agnostic information over instance-specific information.

%\paragraph{Role of critic on data collection and exploration.}
%[WRITE BASED ON FOLLOWING STRUCTURE BELOW]
% - By applying our MI-metrics to the raw collected observations collected by the actor during training, We find that the critic influences data collection. This occurs even when policy and critic networks are fully decoupled. Through its value estimates being used in the policy-gradient objective of the actor, it will co-opt policy updates so that the policy collects data facilitating the optimisation of the critic's own learning objective.
% - This hints as the critic playing an important role in exploration through this mechanism (whic is potentially good) and/or additional interference it exerts on learning the optimal policy, even post decoupling (which is bad).

%\paragraph{Training Instance Augmentation via a learned generative model.}
%[WRITE BASED ON FOLLOWING STRUCTURE BELOW]
% - We introduce Data-Regularised Environment Design (DRED), a method for improving zero-shot generalisation in RL when only a limited set of training instances is initially available.
% - DRED trains a generative model to approximate the distribution represented by the available environment parameters, then samples new instances from that learned distribution via an adaptive sampling strategy to augment agent training.
% - We demonstrate that DRED reduces [I(Z;L)] more effectively than methods sampling from a fixed set of instances, even adaptively. This is turns reduces the generalisation error.

%\paragraph{Importance of distributional grounding for effective zero-shot transfer.}
%[WRITE BASED ON FOLLOWING STRUCTURE BELOW]
% - demonstrate that existing unsupervised environment-design methods can induce substantial distribution shift and actually reduce zero-shot transfer on the target distribution.
% - DRED prevents this failure by constraining generation through a learned approximation to the ground-truth instance distribution. Across the evaluated domains, this grounded augmentation improves zero-shot transfer over prior unsupervised environment-design approaches
% - DRED lets RL agents extracts better zero-shot transfer capbilities from a restricted set of training instances. Thereby reducing the manual scenario-design burden associated with the second condition of the simulation bottleneck.

\section{Thesis Outline}
\label{sec:intro_outline}

\paragraph{Chapter~2} studies the interplay between actor and critic representations in deep RL. It first motivates the distinction between information required for policy optimisation and information required for value estimation. The chapter then introduces mutual-information-based metrics for measuring instance dependence and representational specialisation with respect to actions, values, and latent dynamics. Using these metrics, it compares architectures in which the actor and critic share a feature extractor with architectures in which their representations are learned separately. The results show that separated representations develop complementary specialisations and improve generalisation. The chapter establishes a bound relating the proposed information quantities to generalisation error, evaluates their empirical relationship with zero-shot transfer, and analyses the effects of auxiliary objectives. Finally, we examine the role the critic plays on the interaction data being generated by the learned policy, and we conclude the critic partially co-opts the policy optimisation process to be advantageous to the critic's own learning objective.

\paragraph{Chapter~3} introduces Data-Regularised Environment Design (DRED). The first part of the chapter analyses adaptive level sampling through the representational framework developed in Chapter~2, showing how value-loss prioritisation can regularise shared actor--critic representations. It then studies unsupervised environment-design methods and demonstrates that unconstrained generation can shift the training distribution away from the distribution of interest. The chapter presents DRED, which learns a generative model of environment parameters from an initial dataset and uses this model to generate new, distributionally grounded training instances. Experiments compare DRED with uniform sampling, adaptive sampling, and prior environment-design methods. The results show that DRED expands training diversity while preserving semantic alignment with the target distribution, leading to stronger zero-shot generalisation.

\paragraph{Chapter~4} moves from the generation of instance parameters within a given simulator to the problem of learning a neural simulator from data. It introduces PERSIST, a world model with a dynamic latent 3D representation of the environment. The chapter explains the limitations of relying on bounded pixel and action histories as the sole conditioning mechanism for action-conditioned video generation. Instead, it proposes a decomposition into models predicting the 3D environment state, the camera state, and visual observations. It describes how a latent 3D representation can be updated from interaction and projected into screen space to condition visual generation. Quantitative experiments and a user study assess spatial memory, geometric consistency, visual quality, and long-horizon stability of PERSIST against approaches relying on pixel histories.

\paragraph{Chapter~5} summarizes our main findings, discusses limitations of our work, and points out directions for future research.

\section{Published Work}
\label{sec:intro_published_work}

Portions of this thesis have been published previously. The analysis of actor and critic representation sharing, representational specialisation, auxiliary objectives, and the role of the critic in exploration presented in Chapter~2 appears in \emph{Studying the Interplay Between the Actor and Critic Representations in Reinforcement Learning} \citep{garcin2025studying}. The work on adaptive level sampling, distribution shift in unsupervised environment design, and Data-Regularised Environment Design presented in Chapter~3 appears in \emph{DRED: Zero-shot Transfer in Reinforcement Learning via Data-Regularised Environment Design} \citep{garcin2024dred}. The persistent 3D world-modelling framework presented in Chapter~4 appears in \emph{Beyond Pixel Histories: World Models with Persistent 3D State} \citep{garcin2026beyond}.
